\documentclass[a4paper,11pt]{article}
\usepackage[utf8]{inputenc}
\usepackage[french]{babel}
\usepackage[T1]{fontenc}
\usepackage{amsmath,amssymb,amsthm}
\usepackage{geometry}
\geometry{margin=2cm}
\usepackage{enumitem}
\usepackage{hyperref}
\usepackage{booktabs}
\usepackage{xcolor}
\usepackage{listings}
\usepackage{titlesec}
\usepackage{fancyhdr}
\usepackage{lastpage}
\AtBeginDocument{\shorthandoff{;:!?}}

\titleformat{\section}{\large\bfseries\color{blue!60!black}}{}{0pt}{}[\vspace{-0.5ex}\rule{\textwidth}{0.4pt}]
\titleformat{\subsection}{\bfseries\color{blue!40!black}}{}{0pt}{}

\newtheorem*{important}{\`A retenir}
\pagestyle{fancy}
\fancyhf{}
\fancyhead[L]{\small STA211 -- Notebook Correction (Keras)}
\fancyhead[R]{\small Fiche r\'ecapitulative}
\fancyfoot[C]{\thepage/\pageref{LastPage}}
\renewcommand{\headrulewidth}{0.4pt}

\title{\textbf{Fiche r\'ecapitulative -- STA211}\\[4pt]
\large Notebook Jupyter -- Correction \\[2pt]
\small DL-STA211-correction.ipynb -- Solutions et r\'esultats}
\author{N. Thome (CNAM)}
\date{14 juin 2026}

\begin{document}
\maketitle
\thispagestyle{empty}

\begin{abstract}
Ce notebook ex\'ecut\'e fournit les solutions du TP Keras.
Les performances obtenues et les visualisations t-SNE du manifold
untangling y sont document\'ees.
\end{abstract}

\vfill
\tableofcontents
\newpage

% ===================================================================
\section{R\'esultats des exercices}
% ===================================================================

\subsection{Ex0 -- Chargement}

Ex\'ecution r\'eussie (2 cellules). Donn\'ees MNIST charg\'ees,
reshape, normalisation, affichage des 200 premi\`eres images.

\subsection{Ex1 -- R\'egression logistique}

5 cellules ex\'ecut\'ees.

\begin{itemize}
  \item Mod\`ele~: 7\,850 param\`etres (Softmax \`a 10 classes).
  \item Apprentissage~: SGD, lr=0.5, batch=100, 10 \'epoques.
  \item \textbf{Accuracy test~: $\sim$92\,\%} (confirm\'e).
\end{itemize}

\begin{lstlisting}
# Cell 7 - fit
batch_size = 100
nb_epoch = 10
model.fit(X_train, Y_train, batch_size=batch_size,
          epochs=nb_epoch, verbose=1)
# Cell 8 - evaluate
scores = model.evaluate(X_test, Y_test, verbose=0)
print("%s: %.2f%%" % (model.metrics_names[0], scores[0]*100))
print("%s: %.2f%%" % (model.metrics_names[1], scores[1]*100))
\end{lstlisting}

\subsection{Ex2 -- MLP}

2 cellules ex\'ecut\'ees.

\begin{itemize}
  \item Mod\`ele~: 79\,510 param\`etres (couche cach\'ee 100D).
  \item \textbf{Accuracy test~: $\sim$98\,\%} (sup\'erieur \`a LR).
\end{itemize}

\subsection{Ex3 -- CNN}

3 cellules ex\'ecut\'ees.

\begin{itemize}
  \item Architecture LeNet-like~: Conv16 $\rightarrow$ Pool
    $\rightarrow$ Conv32 $\rightarrow$ Pool $\rightarrow$ Flat
    $\rightarrow$ Dense100 $\rightarrow$ Dense10.
  \item 20 \'epoques.
  \item \textbf{Accuracy test~: $\sim$99\,\%} (meilleur score).
\end{itemize}

\subsection{Ex4 -- t-SNE et m\'etriques de s\'eparation}

12 cellules ex\'ecut\'ees.

\begin{itemize}
  \item t-SNE~: 2 composantes, init='pca', perplexity=30,
    verbose=2.
  \item ConvexHull~: enveloppe convexe par classe.
  \item Ellipses~: GaussianMixture (1 composante,
    covariance\_type='full').
  \item Neighborhood Hit (NH)~: k=6, score moyen sur les 10
    classes.
\end{itemize}

\paragraph{Observation~:} la s\'eparation des classes dans
l'espace t-SNE est nette, contrairement \`a l'ACP (classes
emm\^el\'ees). t-SNE s\'epare bien les vari\'et\'es non
lin\'eaires.

\subsection{Ex5 -- Manifold Untangling}

5 cellules ex\'ecut\'ees.

\begin{itemize}
  \item Chargement du MLP (ou CNN) sauvegard\'e.
  \item Suppression des couches hautes par \texttt{pop()}.
  \item Extraction des repr\'esentations internes (latentes)
    par \texttt{model.predict(X\_test)}.
  \item Projection t-SNE des repr\'esentations.
\end{itemize}

\paragraph{R\'esultat~:} les classes sont bien mieux s\'epar\'ees
dans l'espace latent que dans l'espace d'entr\'ee original.
Le CNN offre une s\'eparation encore meilleure que le MLP.

% ===================================================================
\section{Performances r\'ecapitulatives}
% ===================================================================

\begin{center}
\begin{tabular}{lcc}
\toprule
Mod\`ele & Accuracy test & Temps (20 \'epoques) \\
\midrule
R\'egression logistique & 92\,\% & rapide (CPU) \\
MLP (1 couche cach\'ee) & 98\,\% & mod\'er\'e \\
CNN (LeNet-like) & 99\,\% & lent (GPU recommand\'e) \\
\bottomrule
\end{tabular}
\end{center}

\begin{important}
  La progression des performances refl\`ete la capacit\'e
  croissante des mod\`eles \`a apprendre des repr\'esentations
  de plus en plus abstraites~: lin\'eaire (LR) $\rightarrow$
  non lin\'eaire peu profond (MLP) $\rightarrow$ non lin\'eaire
  profond avec structure spatiale (CNN).
\end{important}

\newpage
\section*{Questions cl\'es (QCM)}
\addcontentsline{toc}{section}{Questions cl\'es (QCM)}

\begin{enumerate}
  \item Accuracy de la r\'egression logistique sur MNIST~?
    \textbf{R. $\sim$92\,\%.}
  \item Accuracy du MLP sur MNIST~? \textbf{R. $\sim$98\,\%.}
  \item Accuracy du CNN sur MNIST~? \textbf{R. $\sim$99\,\%.}
  \item Quel mod\`ele donne le meilleur score~?
    \textbf{R. Le CNN.}
  \item Combien de couches de convolution dans le CNN du TP~?
    \textbf{R. 2 (Conv16 + Conv32).}
  \item Combien de neurones dans la couche cach\'ee du CNN~?
    \textbf{R. 100.}
  \item Quelle m\'ethode de r\'eduction non lin\'eaire est
    utilis\'ee~? \textbf{R. t-SNE.}
  \item Les classes sont-elles mieux s\'epar\'ees dans l'espace
    latent ou dans l'espace d'entr\'ee~?
    \textbf{R. Dans l'espace latent (manifold untangling).}
  \item Quelle m\'ethode Keras supprime la derni\`ere couche~?
    \textbf{R. model.pop().}
  \item Quel optimiseur est utilis\'e pour tous les mod\`eles~?
    \textbf{R. SGD (learning\_rate=0.5).}
\end{enumerate}

\end{document}
